% Chapter 2: Background and Scope
\chapter{法规背景与适用范围}

\section{GB 4785-2019 概况}

\subsection{基本信息}

GB 4785-2019《汽车及挂车外部照明和光信号装置的安装规定》是中华人民共和国国家标准，于 2019 年由国家市场监督管理总局、国家标准化管理委员会联合发布，实施日期为 2020 年 1 月 1 日。

\begin{table}[H]
\centering
\caption{GB 4785-2019 基本信息}
\begin{tabularx}{\textwidth}{|l|X|}
\hline
\textbf{标准编号} & GB 4785-2019 \\
\hline
\textbf{标准名称} & 汽车及挂车外部照明和光信号装置的安装规定 \\
\hline
\textbf{英文名称} & Prescription for installation of external lighting and light-signalling devices \\
\hline
\textbf{发布机构} & 国家市场监督管理总局、国家标准化管理委员会 \\
\hline
\textbf{实施日期} & 2020 年 1 月 1 日 \\
\hline
\textbf{替代标准} & GB 4785-2007 \\
\hline
\textbf{采用关系} & 修改采用 UN R48 系列修订（参考 05、06 系列修订） \\
\hline
\textbf{归口单位} & 全国汽车标准化技术委员会（SAC/TC 114） \\
\hline
\textbf{起草单位} & 中国汽车技术研究中心、清华大学、星宇车灯等 \\
\hline
\textbf{标准属性} & 强制性国家标准 \\
\hline
\textbf{页数规模} & 约 150 页（含附录） \\
\hline
\end{tabularx}
\end{table}

\subsection{修订背景}

GB 4785-2007 实施期间，UN R48 经历了多次修订，特别是 04、05、06 系列修订，引入了大量新技术要求（如 AFS、DRL 等）。为了与时俱进，全国汽车标准化技术委员会启动了修订工作。

GB 4785-2019 的主要修订内容包括：

\begin{itemize}[leftmargin=2em]
    \item 调整范围分类（M 类、N 类、O 类车辆）；
    \item 整合 AFS 相关要求；
    \item 强化昼间行驶灯（DRL）要求；
    \item 调整部分色度坐标；
    \item 修订测试方法；
    \item 增加对新型光源（如 LED）的考量。
\end{itemize}

\subsection{适用范围}

GB 4785-2019 适用于以下车辆：

\begin{itemize}[itemsep=0.2em]
    \item M 类车辆：至少有 4 个车轮的载客车辆，至少有 3 个车轮且最大设计车速超过 25 km/h 的载客车辆；
    \item N 类车辆：至少有 4 个车轮的载货车辆，至少有 3 个车轮且最大设计车速超过 25 km/h 的载货车辆；
    \item O 类车辆：挂车（包括半挂车）。
\end{itemize}

不包括：

\begin{itemize}[itemsep=0.2em]
    \item 设计车速低于 25 km/h 的低速车辆；
    \item 特殊用途车辆（如军用车、消防车）的特殊设备；
    \item 非道路行驶的工程机械。
\end{itemize}

\subsection{装置类型}

GB 4785-2019 涵盖了以下类型的外部照明和光信号装置：

\begin{enumerate}[leftmargin=2em]
    \item 远光灯（Driving beam headlamp, DRBH）；
    \item 近光灯（Passing beam headlamp, PBH）；
    \item 前雾灯（Front fog lamp, FFL）；
    \item 倒车灯（Reversing lamp, RL）；
    \item 转向信号灯（Direction indicator lamp, DIL）；
    \item 危险警告信号（Hazard warning signal）；
    \item 制动灯（Stop lamp, SL）；
    \item 牌照灯（Licence plate lamp, LPL）；
    \item 前位置灯（Front position lamp, FPL）；
    \item 后位置灯（Rear position lamp, RPL）；
    \item 后雾灯（Rear fog lamp, RFL）；
    \item 驻车灯（Parking lamp, PL）；
    \item 昼间行驶灯（Daytime running lamp, DRL）；
    \item 侧标志灯（Side marker lamp, SML）；
    \item 前示廓灯（Front end-outline marker lamp）；
    \item 后示廓灯（Rear end-outline marker lamp）；
    \item 侧反射器（Side retro-reflector）；
    \item 后反射器（Rear retro-reflector）；
    \item 前反射器（Front retro-reflector）；
    \item 自适应前照灯系统（AFS）；
    \item 紧急制动信号（Emergency stop signal, ESS）；
    \item 角灯（Cornering lamp, CL）。
\end{enumerate}

\section{GB 4785 草案 V4.3.5 概况}

\subsection{基本信息}

\begin{table}[H]
\centering
\caption{GB 4785 草案 V4.3.5 基本信息}
\begin{tabularx}{\textwidth}{|l|X|}
\hline
\textbf{文件代号} & GB 4785 草案 V4.3.5 \\
\hline
\textbf{草案名称} & 汽车及挂车外部照明和光信号装置的安装规定（征求意见稿） \\
\hline
\textbf{草案版本} & V4.3.5（征求意见稿 1315 版） \\
\hline
\textbf{编制日期} & 2026 年 8 月 \\
\hline
\textbf{替代标准} & 拟替代 GB 4785-2019 \\
\hline
\textbf{对齐目标} & 主要对齐 UN R48-09，并兼顾中国国情 \\
\hline
\textbf{工作状态} & 征求意见阶段 \\
\hline
\textbf{预期实施} & 2027--2028 年（具体时间待定） \\
\hline
\end{tabularx}
\end{table}

\subsection{修订动因}

启动本次修订的主要动因包括：

\begin{enumerate}[leftmargin=2em]
    \item \textbf{国际对接}：UN R48 已演进至 09 系列修订（2025 年 9 月生效），现行 GB 4785-2019 主要对应 05--06 系列修订，与最新国际要求差距加大；
    \item \textbf{新技术}：ADB、矩阵式 LED 等智能照明技术快速发展，亟需在法规层面作出规定；
    \item \textbf{电动汽车}：新能源汽车（特别是纯电动车）的特殊需求，如低电压条件、能量回收等需要专门考量；
    \item \textbf{智能网联}：V2X（车联网）的发展使得灯具与车辆其他系统的协同更加紧密；
    \item \textbf{事故数据}：近年的交通事故数据分析表明，部分照明场景需要更严格的要求。
\end{enumerate}

\subsection{适用范围}

草案版的适用范围在现行版基础上进行了微调：

\begin{itemize}[leftmargin=2em]
    \item 维持 M、N、O 类车辆的基础分类不变；
    \item \textbf{新增}：明确包含最大设计车速超过 25 km/h 的 L 类车辆中的部分车辆；
    \item \textbf{新增}：明确包含低速电动车（按 GB/T 24158 标准定义的车辆）；
    \item \textbf{明确}：明确将自动驾驶车辆的特殊照明需求纳入考量范围；
    \item \textbf{保留}：继续保留对挂车、半挂车的相关规定。
\end{itemize}

\subsection{主要变化预期}

草案版预期在以下方面产生重大变化：

\begin{enumerate}[leftmargin=2em]
    \item 增加 ADB（自适应远光灯）的安装规定；
    \item 强化 AFS（自适应前照灯系统）的等级转换与故障指示；
    \item 调整 DRL（昼间行驶灯）的安装要求；
    \item 新增或强化 ESS（紧急制动信号）的要求；
    \item 完善对 LED 矩阵式灯具的测试方法；
    \item 修订色度坐标以匹配国际最新数据；
    \item 强化电气安全与电磁兼容性要求；
    \item 增加对自动驾驶车辆的照明协同要求；
    \item 完善对挂车/半挂车与牵引车的联动要求。
\end{enumerate}

\section{UN R48-09 概况}

\subsection{基本信息}

\begin{table}[H]
\centering
\caption{UN R48-09 基本信息}
\begin{tabularx}{\textwidth}{|l|X|}
\hline
\textbf{法规编号} & UN Regulation No. 48 \\
\hline
\textbf{法规名称} & Uniform provisions concerning the approval of vehicles with respect to the installation of lighting and light-signalling devices \\
\hline
\textbf{中文译名} & 关于就照明和光信号装置的安装批准车辆的统一规定 \\
\hline
\textbf{当前修订} & 09 Series of Amendments \\
\hline
\textbf{生效日期} & 2025 年 9 月 26 日 \\
\hline
\textbf{制定机构} & UNECE / WP.29 / GRE \\
\hline
\textbf{法规性质} & 1958 年协议下的 UNECE 法规 \\
\hline
\textbf{适用范围} & M、N、O 类车辆（与 GB 4785 类似） \\
\hline
\textbf{页数规模} & 约 250 页（含附录） \\
\hline
\end{tabularx}
\end{table}

\subsection{修订历程}

UN R48 自 1965 年发布以来，已经历多次重大修订：

\begin{table}[H]
\centering
\caption{UN R48 主要修订历程}
\begin{tabular}{|l|c|p{8cm}|}
\hline
\textbf{修订系列} & \textbf{年份} & \textbf{主要修订内容} \\
\hline
原始版本    & 1965  & 首次发布 \\
\hline
01 系列     & 1993  & 整合更新 \\
\hline
02 系列     & 2001  & 增加新型信号装置 \\
\hline
03 系列     & 2007  & 引入 DRL、AFS 概念 \\
\hline
04 系列     & 2010  & 强化 DRL 和制动灯要求 \\
\hline
05 系列     & 2011  & 引入 ESS、互连系统 \\
\hline
06 系列     & 2012  & 大量技术细节修订 \\
\hline
07 系列     & 2014  & 增加 ADB 相关条款 \\
\hline
08 系列     & 2016  & 引入新光源分类 \\
\hline
09 系列     & 2019--2025 & 全面修订，引入矩阵式 LED 等新技术 \\
\hline
\end{tabular}
\end{table}

\subsection{09 系列修订的主要技术变化}

UN R48-09 相对前序版本的主要技术变化包括：

\begin{itemize}[itemsep=0.3em]
    \item 全面引入 ADB（自适应远光灯）的安装规定；
    \item 完善 AFS 系统的等级转换要求；
    \item 强化 DRL 的光度与色度要求；
    \item 引入新型色度坐标系统；
    \item 引入 LED 光源的特殊考量；
    \item 强化电磁兼容性（EMC）要求；
    \item 引入网络安全基础要求；
    \item 修订测试方法（如暗室测试、路面测试）；
    \item 引入数字化文档与认证。
\end{itemize}

\subsection{适用范围与车辆分类}

UN R48-09 适用于以下车辆类别（与 GB 4785 兼容）：

\begin{itemize}[itemsep=0.2em]
    \item M1 类：载客车辆（$\le$ 9 座，含驾驶员）；
    \item M2、M3 类：载客车辆（$>$ 9 座）；
    \item N1 类：载货车辆（最大总质量 $\le$ 3.5 t）；
    \item N2、N3 类：载货车辆（最大总质量 $>$ 3.5 t）；
    \item O1--O4 类：挂车（按总质量分级）。
\end{itemize}

\section{三版本适用范围对比}

\subsection{车辆类别覆盖范围对比}

\begin{table}[H]
\centering
\caption{三版本适用车辆类别对比}
\begin{tabular}{|l|c|c|c|}
\hline
\textbf{车辆类别} & \textbf{\versionAShort} & \textbf{\versionBShort} & \textbf{\versionCShort} \\
\hline
M 类（载客）           & 是 & 是 & 是 \\
\hline
N 类（载货）           & 是 & 是 & 是 \\
\hline
O 类（挂车）           & 是 & 是 & 是 \\
\hline
L 类（摩托，部分）     & 豁免 & 部分纳入 & 豁免 \\
\hline
低速电动车            & 适用 & 新增明确 & 适用 \\
\hline
自动驾驶车辆         & 一般规定 & 新增明确 & 一般规定 \\
\hline
特殊用途车           & 部分豁免 & 部分豁免 & 部分豁免 \\
\hline
军用车辆              & 豁免 & 豁免 & 豁免 \\
\hline
\end{tabular}
\end{table}

\subsection{装置类型覆盖范围对比}

\begin{table}[H]
\centering
\caption{三版本适用装置类型对比}
\begin{tabular}{|l|c|c|c|}
\hline
\textbf{装置类型} & \textbf{\versionAShort} & \textbf{\versionBShort} & \textbf{\versionCShort} \\
\hline
远光灯                 & 是 & 是 & 是 \\
\hline
近光灯                 & 是 & 是 & 是 \\
\hline
前雾灯                 & 是 & 是 & 是 \\
\hline
后雾灯                 & 是 & 是 & 是 \\
\hline
转向信号灯             & 是 & 是 & 是 \\
\hline
危险警告信号            & 是 & 是 & 是 \\
\hline
制动灯                 & 是 & 是 & 是 \\
\hline
ESS（紧急制动）        & 部分引入 & 强化 & 是 \\
\hline
前/后位置灯            & 是 & 是 & 是 \\
\hline
牌照灯                 & 是 & 是 & 是 \\
\hline
驻车灯                 & 是 & 是 & 是 \\
\hline
DRL（昼间行驶灯）      & 是 & 是 & 是 \\
\hline
侧标志灯               & 是 & 是 & 是 \\
\hline
前/后示廓灯           & 是 & 是 & 是 \\
\hline
前/侧/后反射器        & 是 & 是 & 是 \\
\hline
倒车灯                 & 是 & 是 & 是 \\
\hline
AFS                    & 是 & 强化 & 是 \\
\hline
ADB                    & 未规定 & 新增 & 是 \\
\hline
角灯                   & 是 & 是 & 是 \\
\hline
Logo 投影灯            & 未规定 & 新增 & 是 \\
\hline
迎宾灯                 & 未规定 & 新增 & 部分规定 \\
\hline
\end{tabular}
\end{table}

\subsection{对齐度评估}

基于上述对比，可以初步评估三版本的对齐度：

\begin{table}[H]
\centering
\caption{三版本技术对齐度评估}
\begin{tabular}{|l|c|c|c|}
\hline
\textbf{对比维度} & \textbf{\versionAShort 对 \versionCShort} & \textbf{\versionBShort 对 \versionCShort} & \textbf{\versionBShort 对 \versionAShort} \\
\hline
车辆类别覆盖     & 95\%    & 98\%    & 改进 \\
\hline
装置类型覆盖     & 75\%    & 95\%    & 改进 \\
\hline
技术要求细则     & 70\%    & 88\%    & 改进 \\
\hline
测试方法         & 65\%    & 85\%    & 改进 \\
\hline
术语定义         & 80\%    & 92\%    & 改进 \\
\hline
认证程序         & 60\%    & 82\%    & 改进 \\
\hline
\textbf{综合对齐度} & \textbf{约 70\%} & \textbf{约 90\%} & \textbf{显著提升} \\
\hline
\end{tabular}
\end{table}

\section{小结}

通过对三版本基本信息和适用范围的分析，可以得出以下结论：

\begin{enumerate}[leftmargin=2em]
    \item 草案版相对现行版在适用范围上进行了适度扩展，特别是对新型车辆（低速电动车、自动驾驶车辆）作出明确规定；
    \item 草案版在装置类型覆盖上向 UN R48-09 显著靠拢，新增了 ADB、Logo 投影灯、迎宾灯等新型装置；
    \item 草案版对现行版的核心装置（远光灯、近光灯、信号灯等）均予以保留并强化要求；
    \item 综合对齐度从约 70\% 提升至约 90\%，表明中国标准正在加快与国际接轨的步伐。
\end{enumerate}